% !TeX program = xelatex
\documentclass[11pt,a4paper]{article}
% The article layout follows paper/main.tex in the source preprint.
% XeLaTeX replaces inputenc/fontenc so Russian text is portable to Overleaf.
\usepackage{fontspec,polyglossia}
\setmainlanguage{russian}\setotherlanguage{english}
\setmainfont{DejaVu Serif}[ItalicFont=DejaVu Serif,ItalicFeatures={FakeSlant=0.16},BoldItalicFont=DejaVu Serif Bold,BoldItalicFeatures={FakeSlant=0.16}]
\setsansfont{DejaVu Sans}\setmonofont{DejaVu Sans Mono}[Scale=0.82]
\usepackage{amsmath,amssymb,amsthm,mathtools}
\usepackage{unicode-math}\setmathfont{Latin Modern Math}
\usepackage[margin=24mm,headheight=14pt]{geometry}
\usepackage{microtype,xcolor,booktabs,tabularx,longtable,array,enumitem,fancyhdr,needspace,xurl}
\usepackage[unicode,colorlinks=true,linkcolor=blue!45!black,urlcolor=blue!45!black,citecolor=blue!45!black]{hyperref}
\hypersetup{pdftitle={Выпрямления гармонических вложений, LND и инвариант Макар-Лиманова: исследовательский аппендикс},pdfauthor={Ivan Dorofeev and Alexander Frolov},pdfsubject={Research appendix under verification; 28 October 2026}}
\pagestyle{fancy}\fancyhf{}
\fancyhead[L]{\small Исследовательский аппендикс: LND и ML}
\fancyhead[R]{\small I. Dorofeev and A. Frolov}\fancyfoot[C]{\thepage}
\renewcommand{\headrulewidth}{0.3pt}
\setlength{\parindent}{0pt}\setlength{\parskip}{5pt}\setlength{\emergencystretch}{3em}
\setlist{nosep,leftmargin=*}\setcounter{tocdepth}{1}
\newtheorem{theorem}{Теорема}[section]
\newtheorem{proposition}[theorem]{Предложение}
\newtheorem{corollary}[theorem]{Следствие}
\theoremstyle{definition}\newtheorem{definition}[theorem]{Определение}
\theoremstyle{remark}\newtheorem{remark}[theorem]{Замечание}
\newcommand{\A}{{\mathbb{A}}}\newcommand{\Q}{{\mathbb{Q}}}\newcommand{\C}{{\mathbb{C}}}\newcommand{\R}{{\mathbb{R}}}
\newcommand{\Aut}{\operatorname{Aut}}\newcommand{\ML}{\operatorname{ML}}\newcommand{\LND}{\operatorname{LND}}
\newcommand{\pl}{\operatorname{plinth}}\newcommand{\gr}{\operatorname{gr}}\newcommand{\wt}{\operatorname{wt}}\newcommand{\Div}{\operatorname{div}}\newcommand{\Span}{\operatorname{span}}\newcommand{\pr}{\operatorname{pr}}
\newcommand{\floor}[1]{\left\lfloor#1\right\rfloor}\newcommand{\file}[1]{\path{#1}}
\newcolumntype{L}[1]{>{\raggedright\arraybackslash}p{#1}}
\widowpenalty=10000\clubpenalty=10000
% Reserve enough space for two-digit section numbers in the contents.
\makeatletter
\renewcommand*\l@section[2]{\ifnum\c@tocdepth>\z@\addpenalty\@secpenalty\addvspace{0.8em}\begingroup\parindent\z@\rightskip\@pnumwidth\parfillskip-\@pnumwidth\leavevmode\bfseries\advance\leftskip3em\hskip-\leftskip\@tempdima3em #1\nobreak\hfil\nobreak\hb@xt@\@pnumwidth{\hss #2}\par\endgroup\fi}
\makeatother

\title{\vspace{-12mm}\LARGE\bfseries Выпрямления гармонических вложений, LND\newline и инвариант Макар-Лиманова\\[5pt]\large Исследовательский аппендикс: результаты,\newline точные сертификаты и открытые задачи}
\author{Ivan Dorofeev\textsuperscript{1} and Alexander Frolov\textsuperscript{2}\\[3pt]
\normalsize\textsuperscript{1}Independent Researcher\enspace\textperiodcentered\enspace\href{mailto:qft98romanov@gmail.com}{\texttt{qft98romanov@gmail.com}}\\
\normalsize\textsuperscript{2}Higher School of Modern Mathematics, MIPT}
\date{28 октября 2026 года}
\begin{document}\raggedright
\maketitle\thispagestyle{empty}\vspace{-12mm}
\begin{abstract}
Систематизируются результаты исследовательской ветки, исходящей из явных
комплексных выпрямлений трилистника Шастри и гармонической восьмёрки.
Отделены следствия для конкретных гипотез van den Essen--van Rossum от
общих нерешённых здесь задач вложения, сокращения и линеаризации.
Приведены конструктивные переходы к локально нильпотентным дифференцированиям,
точные минимумы степеней срезов, формальные координаты ядер и формулы
фильтраций для классов аффинных модификаций. Выпрямление прямой
характеризуется существованием нормального локально конечного поля;
даны конечные сертификаты для поиска таких полей. Для $H(3,7,8)$ и
$H(4,5,7)$ найдены полные линейные пространства кандидатов размерностей
$24$ и $19$, но их локальная конечность и новые комплексные выпрямления
не установлены. Приложение содержит воспроизводимые вычисления и точные
границы всех утверждений. Мотивные методы, $K$-теория и гомологическая
алгебра не используются.
\end{abstract}
\smallskip
\textbf{Благодарности.} Авторы выражают искреннюю благодарность
профессору Алексею Канель-Белову и профессору Александру Гасникову.
\begin{quote}\small
\textbf{Исследовательская ветка. Всё представленное в этом аппендиксе
находится на стадии проверки.} Ветка продолжает изучение следствий
найденных выпрямлений и их связей с проблематикой Абъянкара--Моха,
теорией LND, инвариантом Макар-Лиманова и аффинной геометрией.
Авторы сочли полезным поделиться промежуточными результатами, доказательствами
и вычислениями до завершения проверки. Они надеются, что это направление
вырастет в полноценную исследовательскую программу, и приглашают всех
заинтересованных к проверке, обсуждению и совместному исследованию.
\end{quote}
\smallskip
\textbf{Происхождение и статус.} Материал подготовлен в ходе интерактивной
работы с ИИ; отдельные аргументы и реализации проверялись внутри того же
процесса. Это не внешнее независимое рецензирование. Литературный приоритет
не заявляется. Термин «доказано» ниже означает наличие приведённого
математического доказательства и, где применимо, точного сертификата,
а не завершённую проверку сообществом.

Источник: \href{https://github.com/qft98romanov-hash/abhyankar-moh-trefoil}{репозиторий abhyankar-moh-trefoil},
проверенный коммит \texttt{b954f3609e900}.
Текст является отдельным продолжением аппендикса от 26 октября 2026 года.
\newpage\tableofcontents\newpage
\section{Реестр решённых и нерешённых вопросов}
В этой работе важно различать пять статусов: исходный сертифицированный
результат репозитория; его доказанное следствие; новое утверждение рабочего
аппендикса с доказательством; конечный вычислительный результат; открытый
вопрос. Неудача ограниченного поиска не считается доказательством
невыпрямимости. Несколько эквивалентных формулировок одного кандидата
не считаются решениями нескольких общих проблем.

\begin{longtable}{@{}L{0.24\linewidth}L{0.40\linewidth}L{0.30\linewidth}@{}}
\toprule Вопрос & Точный итог & Граница вывода\\\midrule\endhead
Трилистник Шастри & Предъявлены взаимно обратные выпрямители над
$K=\Q(\sqrt{-3})$. & Опровергается невыпрямимость именно этой комплексной
модели; вещественный узел остаётся нетривиальным.\\[5pt]
Гармоническая восьмёрка & Выпрямлено $H(3,5,7)$ с линейным коэффициентом
$+5u$ во второй компоненте. & Эта модель отличается от Example 2
Шастри: там средняя компонента равна $u^5-5u^3+4u$.\\[5pt]
Conjecture 5.1 в \cite{vdER} & Для указанного там $D$ ядро изоморфно
$\C^{[4]}$, в противоположность гипотезе. & Конкретный кандидат не является
контрпримером к сокращению.\\[5pt]
Conjecture 5.2 в \cite{vdER} & Исходный срез является координатой;
его нулевой слой есть $\A^4_\C$. & Следствие полиномиальности того же ядра.\\[5pt]
Conjecture 5.3 в \cite{vdER} & Связанное расширенное кольцо Риса
полиномиально. & В статье это конкретная гипотеза Асанумы;
другие гипотезы Асанумы не затрагиваются.\\[5pt]
Гармонические семейства & Все допустимые $H(3,4,c)$ и $H(3,5,c)$
выпрямляются над $K$. & Не классификация всех $H(a,b,c)$ и не утверждение
о бесконечном числе новых типов вещественных узлов.\\[5pt]
Степени срезов & Полное конечное линейное сведение, восемь точных минимумов;
вся зависимость от кратности для двух моделей. & Степень зависит от исходных
координат и не является инвариантом сопряжения.\\[5pt]
Нормальное поле & Доказан критерий выпрямления через локальную конечность
и конструктивное восстановление координат. & Существование такого поля
для произвольного вложения ещё не доказано.\\[5pt]
$H(3,7,8)$ и $H(4,5,7)$ & Сертифицированы пространства полей-кандидатов,
барьеры степеней и несколько исключённых классов. & Комплексная выпрямимость,
полиномиальность новых ядер и их полный ML здесь не установлены.\\[5pt]
Общие задачи & Пространственная проблема вложения, все гармонические
тройки, сокращение и общая линеаризация здесь не решены. & Утверждение
«не решено этой работой» не заменяет отдельного обзора всех известных
вариантов этих проблем.\\\bottomrule
\end{longtable}

\paragraph{Историческая точность.}
Плоская теорема Абъянкара--Моха--Судзуки уже является классическим
результатом \cite{AM,Suzuki}. Историческое предположение о существовании
невыпрямляемых пространственных прямых не опровергается выпрямлением
одного или двух кандидатов. Вопросы 1 и 2 статьи Шастри \cite{Shastri}
о существовании полиномиальных представлений узлов с комплексным
вложением решены самим Шастри. Наш вывод относится к конкретному
предположению о его трилистнике, зафиксированному в \cite[с.~1319]{vdER}.

\paragraph{Не заявляемые результаты.}
Не доказаны общая гипотеза Абъянкара--Сатайе, прирученность либо дикость
построенных автоморфизмов, минимальность степеней самих выпрямителей,
общая нижняя граница $6n+1$ для степеней срезов и точная общая сложность
разложения последовательных модификаций. Эти вопросы сохранены в
исследовательской программе, но не включены в список закрытых.

\section{Исходные выпрямления и гармонические семейства}
\subsection{Обозначения и проверенная исходная конструкция}
Поле $k$ имеет характеристику нуль. Полиномы Чебышёва нормированы как
\[
C_j(u)=2T_j(u/2),\qquad C_0=2,\quad C_1=u,\quad C_{j+1}=uC_j-C_{j-1}.
\]
Рассматриваются замкнутые вложения, то есть $k[f_1,f_2,f_3]=k[u]$.
Полином параметра $P$ удовлетворяет $P\circ\gamma=u$; координата, кроме
того, должна дополняться до полиномиального автоморфизма. Эти два
условия не равнозначны.

Исходный препринт \cite{Repository} содержит автоморфизмы
$\Phi_T,\Phi_H$ и их полиномиальные обратные над
$K=\Q[a]/(a^2-a+1)\simeq\Q(\sqrt{-3})$ для
\[
\begin{aligned}
\gamma_T(u)&=(u^3-3u,u^4-4u^2,u^5-10u),\\
\gamma_H(u)&=(C_3(u),C_5(u),C_7(u)).
\end{aligned}
\]
Проверены $\Phi_T\gamma_T=\Phi_H\gamma_H=(u,0,0)$, обе обратные
композиции посредством точных полиномиальных тождеств и все двенадцать
напечатанных компонент. Степени компонент равны соответственно
$(5,15,24)$ и $(7,21,41)$; их минимальность не утверждается.
Полные формулы сохранены в исходном препринте и в автономном
вычислительном пакете. Замена
\[
M(x,y,z)=(x,y+2,z-5x)
\]
переводит $\gamma_T$ в $(C_3,C_4,C_5)$.

\subsection{Конструктивный перенос через модификацию}
Пусть $A=(x,H,W)$ и $A^{-1}=(x,\mathcal Y,\mathcal Z)$ ---
взаимно обратные автоморфизмы, $m\in k[x]$, $\deg m>0$, $p\in k[x,y]$.
Рассмотрим $\beta=(x,p(x,y)+m(x)z,y)$. В факторкольце
$S=k[x]/(m)$ положим $c(y)=H(x,p(x,y),y)$.
Если линейный коэффициент $c$ обратим, а все старшие коэффициенты
нильпотентны, то подстановка $y\mapsto c(y)$ задаёт автоморфизм $S[y]$.
Его композиционный обратный $\psi$ строится
конечным подъёмом по нильпотентному идеалу старших коэффициентов.

Выбираем полиномиальный представитель
\[
\varphi(x,T)\equiv W(x,p(x,\psi(T)),\psi(T))\pmod m.
\]
Тогда
\[
\mathcal A=\left(x,H\circ\beta,
\frac{W\circ\beta-\varphi(x,H\circ\beta)}m\right)
\]
и
\[
\mathcal A^{-1}(x,T,Q)=
\left(x,\mathcal Z(x,T,mQ+\varphi),
\frac{\mathcal Y(x,T,mQ+\varphi)-p(x,\mathcal Z(x,T,mQ+\varphi))}m\right)
\]
имеют полиномиальные компоненты: оба числителя делятся на полный $m$,
включая кратности его корней. После обращения $m$ обе композиции
тождественны. Инъективность локализации переносит равенства в исходное
кольцо, включая исключительный слой $m=0$.
Это конкретная реализация классических механизмов подъёма и относительных
координат; общая идея переноса не объявляется новой \cite{EVM,Repository}.
Случай $m\in k^*$ проще: сама $\beta$ уже является автоморфизмом.

Если $H\circ\gamma=u$ и $w=W\circ\gamma$, то из относительной пары
получается полный выпрямитель
\[
\Phi=(H,\ x-f_1(H),\ W-w(H)).
\]
Наличие только полинома параметра, без проверенной координатной пары,
для этого вывода недостаточно.

\subsection{Полиномиальные изменения гармонических индексов}
Для взаимно простых $a,b$ и знаков $\varepsilon_a,\varepsilon_b\in\{\pm1\}$
условия
\[
c'\equiv\varepsilon_a c\pmod{2a},\qquad
c'\equiv\varepsilon_b c\pmod{2b}
\]
дают
\[
C_{c'}-\varepsilon_a\varepsilon_b C_c\in\Q[C_a,C_b].
\]
Следовательно, соответствующие пространственные кривые связаны
полиномиальным сдвигом третьей координаты.

Краткое доказательство использует узлы плоской проекции
$(C_a,C_b)$. Её кольцо состоит ровно из полиномов по $u$, принимающих
одинаковые значения на двух прообразах каждого узла: включение очевидно,
а равенство следует сравнением коразмерностей
$(a-1)(b-1)/2$ через полугруппу $\langle a,b\rangle$ и интерполяцию.
Указанные сравнения индексов дают требуемые равенства значений по
тригонометрической формуле Чебышёва. Исключение ведущих членов даёт
рациональный полином сдвига. Подробности и точные примеры содержатся
в исходном аппендиксе; литературная основа узловой проекции --- \cite{FF}.

Отсюда следуют два полных семейства над $K$:
\[
H(3,4,c),\quad\gcd(c,12)=1;
\qquad H(3,5,c),\quad\gcd(c,15)=1.
\]
Для пары $(3,4)$ достаточно представителей $c=1,5$ по модулю $24$;
для $(3,5)$ --- $c=1,2,4,7$ по модулю $30$. Представители с индексом
$1$ очевидны, с индексом $2$ известны классически, $c=4$ сводится
перестановкой к $(3,4,5)$, а оставшиеся два обеспечены исходными
выпрямлениями.

Индекс $1$, допустимые модели $H(2,b,c)$ и $H(a,b,a+b)$ уже были
известны в литературе \cite{FF}; они не приписываются этой работе.
Топологическая классификация гармонических узлов \cite{KP} также не
является классификацией комплексных вложений по автоморфизмам.

\subsection{Конечный перебор и его точная область}
Перебраны $44142$ попарно взаимно простые тройки
$2\le a<b<c\le100$ по всем строго убывающим элементарным редукциям
\[
c=ra+sb\quad\longmapsto\quad |ra-sb|,
\qquad (x,y,z)\longmapsto(x,y,C_r(x)C_s(y)-z).
\]
Из них $3209$ достигают индекса $1$, ещё $237$ достигают проверенной
базовой модели, а $40696$ остаются неразрешёнными \emph{этим алгоритмом}.
Последнее число не является числом невыпрямляемых кривых.

Далее для любого вложения $\gamma=(f_1,\ldots,f_n)$ используется
\[
A=k[t,u,X_1,\ldots,X_n],\qquad
D_\mu=t^\mu\partial_u+\sum_i f_i'(u)\partial_{X_i},\qquad
V=\gamma(u)-t^\mu X,
\]
где $\mu\ge1$. Это LND на $n+2$ переменных. Его нельзя смешивать
с якобианным LND модификации на трёх переменных.

% Consolidated proofs from the 26 October continuation; citations normalized.
\section{Теорема о минимальной степени среза}\label{sec:jets}
Положим $V=\gamma(u)-t^\mu X$. Для любого $P\circ\gamma=u$ имеется полиномиальный срез
\[
S_{P,\mu}=\frac{u-P(V)}{t^\mu},\qquad D_\mu(S_{P,\mu})=1.
\]
Полиномиальность следует из точной делимости числителя; знаменатель не означает рациональную функцию. Эта классическая конструкция восходит к van den Essen--van Rossum \cite{vdER}. Наше дальнейшее вычисление уточняет задачу минимизации степени.

\subsection*{Теорема 3.1. Полнота сведения}
Минимум степени среди всех полиномиальных срезов равен минимуму $\deg S_{P,\mu}$ по всем полиномам параметра $P$.

\textit{Доказательство.} Введём веса
$\wt(t)=1$, $\wt(u)=0$, $\wt(X_i)=-\mu$. Дифференцирование повышает вес на $\mu$. Если $D_\mu(s)=1$, компонента $s_{-\mu}$ веса $-\mu$ тоже является срезом; её обычная степень не больше степени $s$.
После обращения $t$:
\[
A_t=k[t,t^{-1},u,V],\qquad D_\mu=t^\mu\partial_u.
\]
Элемент $t^\mu s_{-\mu}-u$ инвариантен и имеет вес нуль. Поскольку все $V_i$ имеют вес нуль, он равен $-P(V)$ для единственного $P\in k[V]$. Редукция точного полиномиального равенства по модулю $t$ даёт $P\circ\gamma=u$.\hfill$\square$

\subsection*{Теорема 3.2. Точная формула и конечный критерий}
С соглашением $\deg0=-\infty$ имеем
\[
\boxed{\deg S_{P,\mu}=\max_{\alpha\ne0}
\bigl\{\deg_u((\partial^\alpha P)\circ\gamma)
+(\mu+1)|\alpha|-\mu\bigr\}.}
\]
Действительно, формула Тейлора даёт
\[
S_{P,\mu}=\sum_{\alpha\ne0}\frac{(-1)^{|\alpha|+1}}{\alpha!}
t^{\mu(|\alpha|-1)}X^\alpha(\partial^\alpha P)(\gamma(u)).
\]
Различным $\alpha$ соответствуют различные мономы по $X$, поэтому сокращения между ними невозможны. Производная порядка $\deg P=r$ даёт нижнюю оценку $(\mu+1)r-\mu$.

Следовательно, срез степени $\le d$ существует тогда и только тогда, когда разрешима конечная линейная система
\[
\begin{gathered}
\deg P\le\floor{\frac{d+\mu}{\mu+1}},\qquad P\circ\gamma=u,\\
[u^j](\partial^\alpha P)(\gamma(u))=0\quad
\text{при }j>d-(\mu+1)|\alpha|+\mu,\quad\alpha\ne0.
\end{gathered}
\]
Она линейна по коэффициентам $P$. Для $n=3$ число неизвестных не превышает
$\binom{\lfloor(d+\mu)/(\mu+1)\rfloor+3}{3}$.
Это полный критерий, а не ограничение поиска выбранным видом среза.

\subsection*{Теорема 3.3. Все достаточно большие кратности}
Пусть
\[
M=\max_i\deg f_i,\qquad
r_\gamma=\min\{\deg P:P\circ\gamma=u\},\qquad
m_\gamma(\mu)=\min_{D_\mu(s)=1}\deg s.
\]
Тогда
\[
\boxed{m_\gamma(\mu)=(\mu+1)r_\gamma-\mu\quad\text{при }\mu\ge M-1.}
\]
Для непостоянного $P$ степени $r$ производная порядка $q$ после подстановки имеет степень не больше $M(r-q)$. При $\mu+1\ge M$:
\[
M(r-q)+(\mu+1)q-\mu\le(\mu+1)r-\mu.
\]
Ненулевая производная порядка $r$ достигает этой границы. Минимизация по $P$ завершает доказательство. Кроме того, $m_\gamma(\mu)$ не убывает: каждое выражение в максимуме имеет неотрицательный наклон $|\alpha|-1$.

Эта метрика зависит от исходных координат. Она не является инвариантом сопряжения LND и не может без дополнительного аргумента служить препятствием полиномиальности ядра.

\section{Точные вычисления для гармонических моделей}\label{sec:computed}
Для $\mu=1$ получены следующие результаты. Нижние границы удостоверены рациональными дуальными сертификатами, верхние --- явными срезами.

\begin{center}\small\renewcommand{\arraystretch}{1.22}
\begin{tabular}{@{}lrrr@{}}\toprule
Модель & $m_\gamma(1)$ & Неизвестных & Ненулевых весов\\\midrule
$H(3,4,5)$ &7&20&11\\
$H(3,5,7)$ &11&56&35\\
$H(4,5,7)$ &12&84&41\\
$H(3,7,8)$ &13&84&41\\
$H(3,10,11)$ &19&220&111\\
$H(3,13,14)$ &25&455&225\\
$H(5,7,8)$ &11&56&29\\
$H(7,10,11)$ &15&120&61\\\bottomrule
\end{tabular}\end{center}
Последние два столбца относятся к системе, исключающей степень на единицу меньше минимума. Первые две строки воспроизводят прежние результаты более компактными сертификатами.

Если система записана $A c=b$, нижний сертификат есть рациональный вектор $\lambda$ с $\lambda^{\mathsf T}A=0$ и $\lambda^{\mathsf T}b=1$. Это исключает решение над любым полем характеристики нуль, а не только над $\Q$. Независимый проверяющий скрипт заново строит столбцы и выполняет вычисления с \file{Fraction}; символьный решатель для проверки не нужен.

\subsection*{$H(3,7,8)$: меньшая степень ретракции не всегда лучше}
Тождества Чебышёва дают
\[
C_7C_8=C_{15}+C_1,\qquad C_5(C_3)=C_{15},\qquad
P_5=yz-C_5(x)=yz-x^5+5x^3-5x.
\]
Полученный срез степени $13$ оптимален при $\mu=1$. Однако минимальная степень самого полинома параметра равна $4$; можно взять
\[
\begin{split}
P_4={}&x^3z+x^3-x^2yz-xy^2-2xz-x\\
&+y^3+2yz-2y.
\end{split}
\]
При $\mu=1$ этот $P_4$ даёт срез степени $19$. Таким образом, минимизировать $\deg P$ и минимизировать $\deg S_P$ --- разные задачи.

Полная зависимость имеет вид
\[
\boxed{\begin{aligned}
m_{378}(1)&=m_{378}(2)=13,\quad m_{378}(3)=17,\quad m_{378}(4)=21,\\
m_{378}(5)&=22,\quad m_{378}(6)=23,\quad
m_{378}(\mu)=3\mu+4\quad(\mu\ge7).
\end{aligned}}
\]
Для $\mu\le4$ достаточно $P_5$, начиная с $\mu=5$ --- $P_4$. Нижние границы при $1\le\mu\le7$ проверены точно; последний тест одновременно исключает $\deg P\le3$. Далее применяется теорема 3.3.

\subsection*{$H(4,5,7)$: формула для всех кратностей}
Здесь
\[
P_3=(x^2-2)z-y^3+3y,\qquad
\boxed{m_{457}(\mu)=\max\{12,\mu+9,2\mu+3\}.}
\]
Первые, вторые и третьи производные $P_3$ вдоль кривой имеют максимальные степени $11$, $7$, $0$. Это даёт верхнюю оценку. При $\mu=1,2,3$ нижняя оценка $12$ следует из сертификата при $\mu=1$ и монотонности.

Ограничение полиномов степени $\le3$ на кривую инъективно: соответствующий минор размера $20\times20$ имеет определитель $267386$. Поэтому $P_3$ единственен в этой степени. При $\mu=4,5$ любой $P$ степени $\ge4$ даёт степень среза хотя бы $3\mu+4>\mu+9$. При $\mu\ge6$ действует общая теорема с $M=7$, $r_\gamma=3$.

\subsection*{Бесконечное семейство верхних границ}
Для нечётных $a,m\ge3$ положим $b=(ma-1)/2$, $c=(ma+1)/2$. Индексы попарно взаимно просты, и $P=yz-C_m(x)$ восстанавливает параметр. Точная степень его среза:
\[
\deg S_{P,\mu}=\max\{a(m-1)+1,(\mu+1)m-\mu\}.
\]
Для $(a,b,c)=(3,3n+1,3n+2)$ получаем $m_\gamma(1)\le6n+1$. Равенство удостоверено при $n=1,2,3,4$; общая нижняя граница пока не доказана.

\section{Ядро и формальная тривиальность}\label{sec:formal}
Пусть $I=\ker(k[Y_1,\ldots,Y_n]\to k[u])$ --- идеал вложения, $B_\mu=\ker D_\mu$. Через разложение Тейлора получается классическое по своей природе описание
\[
\boxed{B_\mu=k[t,V,I(V)/t^\mu]\simeq k[Y][t,I/t^\mu].}
\]
Для $\mu=1$ соответствующая конструкция описана в \cite[Proposition 5.4, Corollary 5.5]{vdER}. Она не утверждает полиномиальности $B_\mu$.

Элементарный ключ: если все производные $F$ порядков $<q$ исчезают вдоль $\gamma$, то $F\in I^q$. Для доказательства выбираем $P\circ\gamma=u$ и $g_i=Y_i-f_i(P(Y))$. Разложение Тейлора вокруг $\gamma(P(Y))$ показывает сначала $I=(g_1,\ldots,g_n)$, затем требуемую принадлежность. Обратное включение следует из правила Лейбница. Весовая компонента $t^jF(V)$ полиномиальна при $j<0$ ровно тогда, когда $F\in I^{\lceil-j/\mu\rceil}$; это даёт формулу ядра.

\subsection*{Теорема 5.1. Все конечные утолщения тривиальны}
Для любого вложения, любого $\mu\ge1$ и $N\ge1$:
\[
\boxed{B_\mu/(t^N)\simeq(k[t]/(t^N))^{[n]}.}
\]
Изоморфизмы можно выбрать совместимыми по $N$, поэтому
\[
\boxed{\widehat B_\mu^{\,(t)}\simeq k[p,q_1,\ldots,q_{n-1}][[t]].}
\]
Справа стоят ряды по $t$ с полиномиальными коэффициентами. Степени коэффициентов по $p,q$ могут расти. Это не утверждение об обычном полиномиальном кольце $k[[t]][p,q]$.

\textit{Доказательство.} Положим $a(u)=\nabla P(\gamma(u))$, $v(u)=\gamma'(u)$. По правилу цепочки $a v=1$. Над PID $k[u]$ модуль строк, аннулирующих $v$, свободен ранга $n-1$; выберем базис $\ell_1,\ldots,\ell_{n-1}$. Вместе с $a$ эти строки образуют обратимую над $k[u]$ матрицу.

Строку $\ell$ можно явно поднять в идеал кривой:
\[
F_\ell(Y)=\sum_i\ell_i(P(Y))\bigl(Y_i-f_i(P(Y))\bigr).
\]
Имеем $F_\ell\in I$ и $\nabla F_\ell(\gamma)=\ell-(\ell v)a=\ell$. Определим
\[
p=P(V),\qquad q_j=-F_{\ell_j}(V)/t^\mu,\qquad s=(u-P(V))/t^\mu.
\]
Это полиномы; $p,q_j$ инвариантны, $D_\mu(s)=1$. При $t=0$ замена $(u,X)\mapsto(p,q,s)$ равна
\[
(u,X)\longmapsto(u,\ell_1(u)X,\ldots,\ell_{n-1}(u)X,a(u)X)
\]
и обратима.

После композиции с её полиномиальной обратной при $t=0$ получаем эндоморфизм $\Phi\equiv\mathrm{id}\pmod{t^\mu}$. На $A/(t^N)$ оператор $E=\mathrm{id}-\Phi$ повышает $t$-адический порядок на $\mu$, поэтому $E^L=0$ при $L=\lceil N/\mu\rceil$. Обратное отображение есть конечная сумма $\sum_{j=0}^{L-1}E^j$. Это аддитивное обратное к гомоморфизму колец, следовательно, оно автоматически мультипликативно.

Итак,
\[
A/(t^N)=(k[t]/(t^N))[p,q_1,\ldots,q_{n-1},s],\qquad\overline D_\mu=\partial_s.
\]
Глобальная теорема о срезе даёт $A=B_\mu[s]$. Поэтому ядро редуцированного дифференцирования равно $B_\mu/(t^N)$. Те же полиномы $p,q$ работают для всех $N$, обеспечивая совместимость.\hfill$\square$

\textbf{Следствие для стратегии.} Абстрактный тип одного конечного утолщения слоя $t=0$ не может отличить эти ядра от полиномиального. Нужно контролировать глобальную алгебраизацию, степени обратного отображения или дополнительные структуры. Конечные сертификаты степенных ограничений по-прежнему полезны: они учитывают исходные координаты, которых нет в абстрактном кольце утолщения.

\section{Ограничение класса координат для кубического семейства}\label{sec:uniform}
Пусть $\gamma_n=(C_3,C_{3n+1},C_{3n+2})$, $n\ge1$. Над $\C$ не существует полиномиальной координаты $H$ с $H\circ\gamma_n=u$ ни одной из форм
\[
\boxed{H=Q(x,y)+M(x)z\quad\text{или}\quad H=Q(x,z)+M(x)y.}
\]
Степени $Q$ и $M$ не ограничиваются. Запрет относится и к абсолютной координате: сохранять $x$ в полном автоморфизме не требуется.

\subsection*{Доказательство относительной части}
Пишем $U_j=U_j(x/2)$ для полиномов Чебышёва второго рода. В свободном модуле с базисом $1,u,u^2-2$ над $\C[x]$, где $u^3-3u=x$, имеем
\[
y=U_nu-U_{n-1}(u^2-2),\qquad
z=-U_{n-1}u+U_n(u^2-2).
\]
Для первой формы редуцируем $Q$ по моническому соотношению
$y^3-3y-C_{3n+1}(x)=0$ до $Ay^2+By+C$. Сравнение коэффициентов даёт
\[
U_{3n}A+U_{2n}M=U_{n-1}.\tag{6.1}
\]
При $M=0$ равенство невозможно по степеням. Если $(x,H,W)$ автоморфизм и $M(\alpha)=0$, полином $Q(\alpha,y)$ должен быть аффинным по $y$. Поэтому $A(\alpha)=0$ и $U_{n-1}(\alpha)=0$.

Все корни $M$ лежат среди вещественных корней $U_{n-1}$. Два наибольших корня $U_{3n}$,
$x_j=2\cos(j\pi/(3n+1))$, $j=1,2$, лежат правее всех этих корней. При $n=1$ корней $U_0$ вообще нет. Подстановка в (6.1) и синусная формула дают
$M(x_1)=1$, $M(x_2)=-1$.
Но для $M=c\prod(x-\alpha)^{e_\alpha}$ отношение $M(x_1)/M(x_2)$ положительно, даже при комплексном $c$. Противоречие.

Для второй формы аналогично получается
\[
U_{3n+1}A+U_{2n}M=U_n.
\]
Два наибольших корня $U_{3n+1}$ лежат правее корней $U_n$ и опять принуждают противоположные значения $M$. Это завершает относительное доказательство.

\subsection*{Переход к абсолютной координате}
При $M\ne0$ полином $H$ является координатой над $\C(x)$. Если он был бы абсолютной координатой, Proposition 3.1 работы Edo--van den Essen--Maubach \cite{EVM} делала бы его координатой над $\C[x]$. Противоречие уже доказанному.

Для $n=1$ запрет совместим с имеющимися в проекте комплексными конструкциями для $H(3,4,5)$: параметрическая компонента должна нелинейно зависеть от обеих поперечных переменных. Это ограничение формы выпрямителя, а не доказательство его отсутствия.

\subsection*{Естественная модификация тоже не проходит}
Рассмотрим
\[
\beta_n=(x,C_n(x)z-y^2+2,y),\qquad
\beta_n\circ\gamma_n=(C_3,C_2,C_{3n+1}).
\]
На каждом корне $\alpha$ многочлена $C_n$ центр $H(\alpha,2-y^2,y)$ не может быть аффинным по $y$, если $H$ восстанавливает параметр целевой тройки. Иначе в факторе по $u^3-3u-\alpha$ получаем
\[
u=\lambda\bigl(U_n(\alpha/2)u-U_{n-1}(\alpha/2)(u^2-2)\bigr)+\nu.
\]
Здесь $U_{n-1}(\alpha/2)\ne0$, поскольку
$C_n(x)^2+(4-x^2)U_{n-1}(x/2)^2=4$. Сравнение коэффициентов сначала даёт $\lambda=0$, затем противоречие. Это доказательство не имеет ограничения степени $H$.

При нечётном $n$ целевая тройка не является вложением: $u=\pm\sqrt3$ имеют одинаковые образы. При чётном $n=2m$ цель имеет явную относительную координату. Если
$C_{6m+1}(u)=uP_m(u^2-2)$, $\varepsilon=P_m(1)=(-1)^m$, то
\[
A_m(Y)=\frac{1-\varepsilon P_m(Y)}{Y-1},\qquad
H_0=xA_m(Y)+\varepsilon Z
\]
восстанавливает параметр. Таким образом, даже выпрямляемость цели не исправляет плохой центр данного переноса.

\subsection*{Дополнительный конечный результат для $H(3,7,8)$}
Исключены координаты $Q(x,L)+M(x)N$, где $(L,N)$ --- любая постоянная невырожденная линейная замена $(y,z)$, при $\deg M\le10$ и произвольной степени $Q$. После нормировки $L=y+az$ задача сводится к 56 точным условиям на кратности корней $M$. Во всех случаях НОД необходимых полиномиальных условий равен $1$; повторные корни при $a=\pm i/2$ проверены отдельно. Это \textbf{ограниченный результат}: степень $M>10$ здесь не исключена.

\section{H(4,5,7): барьер степени четыре в относительном классе}\label{sec:h457}
Пусть $F\circ(C_4,C_5,C_7)=u$ и $\deg F\le4$. Тогда:
\begin{itemize}
\item единственный такой $F$, являющийся координатой $k(x)[y,z]$, есть $F_0=(x^2-2)z-y^3+3y$;
\item ни один такой $F$ не является координатой $k(y)[x,z]$ или $k(z)[x,y]$;
\item $F_0$ не является координатой над $k[x]$: его производные по $y,z$ одновременно исчезают при $x^2=2$, $y^2=1$.
\end{itemize}
Следовательно, для любого метода, дополняющего одну из пар $(x,F)$, $(y,F)$, $(z,F)$ до автоморфизма, обязательно $\deg F\ge5$.

Доказательство исчерпывает всё пространство полиномов степени $\le4$. Из 35 мономов матрица ограничения на кривую имеет ранг 28. Все полиномы параметра образуют аффинное пространство $F_0+\langle K_0,\ldots,K_6\rangle$. В архиве полностью выписаны семь $K_i$; полнота проверена ненулевыми минорами.

Далее применяется теорема Юнга--ван дер Кулка: координата плоскости степени 3 имеет вид $P_3(L)+cM$, а степени 4 --- либо $P_4(L)+cM$, либо $A(M+Q_2(L))^2+bL+c$ с соответствующими ненулевыми коэффициентами. Последовательное сравнение коэффициентов во всех трёх проекциях исключает варианты. В последней ветви вычисляется ненулевой результант $6561301239$; имеется явное тождество Безу с постоянной $729033471$.

Верификатор проверяет 35 точных условий, а отдельный аудит повторяет ранги, граничные случаи и устранение параметров. Структурный источник для нормальных форм --- классическая теорема и её лемма о делимости степеней; удобный первичный текст указан в \cite{Chau}.

\textbf{Граница утверждения.} Оно не исключает абсолютную координату степени 4, завершение которой меняет все три исходные координаты, и не исключает относительные координаты степени $\ge5$.

\section{Полная фильтрация одношагового переноса}\label{sec:transfer}
Пусть $R=k[x]$, $A=R[y,z]$, $(H,W)\in\Aut_R R[Y,Z]$ и
\[
d=H_YW_Z-H_ZW_Y\in k^*,\qquad
\beta=(p(y)+mz,y),\quad p\in R[y],\ 0\ne m\in R.
\]
Положим $h=H\circ\beta$, $w=W\circ\beta$, $D=h_z\partial_y-h_y\partial_z$.
Вынесем содержание по непостоянным мономам:
\[
c=h(x,0,0),\qquad h=c+g h_0,
\]
где $g\in R$ --- монический НОД коэффициентов всех непостоянных по $y,z$ мономов $h$.

\subsection*{Теорема 8.1}
Имеем $g\mid m$, $B=\ker D=k[x,h_0]$. Для каждого неприводимого $\pi\mid m_1=m/g$ определим
\[
E_\pi=v_\pi(m_1),\qquad e_\pi=\deg_y(h_0\bmod\pi)\ge1.
\]
Редукция $h_0$ здесь не зависит от $z$. Для $\mathcal F_n=\ker D^{n+1}$:
\[
\boxed{D^n(\mathcal F_n)=
\left(g^n\prod_{\pi\mid m_1}\pi^{E_\pi(n-\lfloor n/e_\pi\rfloor)}\right)B.}
\]
В частности,
\[
\boxed{\pl(D)=D(A)\cap B=
\left(g\prod_{e_\pi\ge2}\pi^{E_\pi}\right)B.}
\]
Все кратности исключительного делителя сохраняются. Хорошие компоненты $e_\pi=1$ исчезают из plinth-идеала примитивного дифференцирования; множитель $g$ учитывается отдельно.

\subsection*{Основные шаги доказательства}
\textit{1. Ядро и нормировка.} Над $k(x)$ пара $(h,w)$ является системой координат, $D(w)=dm$. Отсюда $g\mid m$ и общее ядро равно $k(x)[h_0]$. Редукция $h_0$ по любому простому $\pi\in R$ непостоянна; поэтому $R[h_0]$ насыщено по $\pi$ в $A$. Последовательное сокращение знаменателей даёт $A\cap k(x)[h_0]=R[h_0]$.

\textit{2. Точное уравнение модификации.} Пусть $(\mathcal Y,\mathcal Z)$ --- обратное к $(H,W)$ отображение. Полином
\[
N(T,V)=\frac{\mathcal Y(c+gT,V)-p(\mathcal Z(c+gT,V))}{g}
\]
действительно лежит в $R[T,V]$, причём $\gcd(N,m_1)=1$; это доказывается редукцией по $g$ и обратимостью центра над конечным факторкольцом. Затем
\[
A\simeq R[T,V,z]/(m_1z-N(T,V)),\qquad T=h_0,\quad V=w.
\]
Полиномиальность $N$ существенна: формальное деление без её проверки оставило бы пробел.

\textit{3. Моническая модель над каждым простым.} Локализуем только $R$: $O=R_{(\pi)}$, коэффициентное кольцо $O[T]$. Над $O/(\pi^{E_\pi})$ центр имеет вид $(G(y),J(y))$, порождающий всё кольцо полиномов по $y$. Оно свободно ранга $e_\pi$ над кольцом по $G$; доказательство получается повторением обычного деления и конечным подъёмом по степеням $\pi$.
Характеристический полином умножения на $J$ даёт монический $P_\pi(T,V)$ степени $e_\pi$ по $V$. Уравнения центра и $P_\pi$ порождают один идеал по модулю $\pi^{E_\pi}$, что приводит к точной модели
\[
P_\pi(h_0,w)=\pi^{E_\pi}v.
\]
Здесь сохраняются условия, зависящие от $T$: замена $O[T]$ полем рациональных функций по $T$ была бы недостаточна.

\textit{4. Базис и старший коэффициент.} Мономы $w^iv^j$, $0\le i<e_\pi$, образуют базис над $O[h_0]$. Их $D$-степени равны $i+e_\pi j$ и попарно различны. Поэтому старшие члены не сокращаются. При $n=i+e_\pi j$ получаем показатель $E_\pi(n-j)$. Пересечение локальных условий по всем простым $R$, с использованием китайской теоремы об остатках для построения предобразов, даёт глобальную формулу.

Полное алгебраическое доказательство, включая делимость, монический подъём и склеивание, приведено в \file{relative_lnd/PROOF.md}; отдельный аудит находится в \file{literature/independent_image_ideal_review.md}. Дополнительно проверены 16 примеров и 28 независимых точных условий.

\subsection*{Пример с кратным делителем и нелинейным подъёмом}
Пусть
\[
h=y^2+xy^3+x^2z,\qquad D=x^2\partial_y-(2y+3xy^2)\partial_z.
\]
Здесь $E=2$, $e=2$. Для $v=y^4+xyz-z$:
\[
Dv=2y+xh,\qquad D^2v=2x^2,\qquad
z=-(1+xy)v+hy^2.
\]
Последовательность идеалов образов имеет показатели $0,2,2,4,4,6,6,\ldots$. Один найденный элемент образа даёт лишь включение; равенство всего идеала следует из теоремы.

\subsection*{Следствия для чистой теории LND}
Для $D_{M,e}=x^M\partial_y-e y^{e-1}\partial_z$, $M\ge1$, $e\ge2$:
\[
I_n(D_{M,e})=(x^{M(n-\lfloor n/e\rfloor)}).
\]
Показатель $M$ восстанавливается из кратности простого делителя plinth-идеала, а $e$ --- как первое $n$ с $I_n\supsetneq I_1^n$. Поэтому разные пары $(M,e)$ дают несопряжённые LND. Само семейство классическое; приоритет этого способа записи следствия не заявляется.

Для $h=(x^2-2)z-y^3+3y$, возникающего у $H(4,5,7)$:
\[
\ker D=k[x,h],\qquad I_n(D)=(x^2-2)^{n-\lfloor n/3\rfloor}k[x,h].
\]
Это точная информация о данном LND, но не препятствие для всех возможных выпрямителей кривой.

\section{Последовательные модификации и группы Кремоны}\label{sec:iterated}
Определим
\[
r_{-1}=z,\quad r_0=y,\quad r_j=xr_{j-2}-r_{j-1}^2,\qquad j\ge1.
\]
Для фиксированного $\ell\ge1$ положим $p=r_\ell$, $w=r_{\ell-1}$ и
\[
D_\ell=(-1)^\ell\bigl(p_y\partial_z-p_z\partial_y\bigr).
\]
Тогда $D_\ell(w)=x^\ell$, $\ker D_\ell=B=k[x,p]$. Дифференцирование неприводимо: общий делитель компонент должен делить $x^\ell$, а $p\bmod x=-y^{2^\ell}$ имеет ненулевую производную по $y$.

На открытом множестве $x\ne0$ каждая операция обратима; знаменатели появляются при продолжении через $x=0$. Это конкретная связь с преобразованиями Кремоны: биррациональная обратимость проста, а полиномиальность и длина разложения контролируются исключительным делителем.

\subsection*{Теорема 9.1. Вся фильтрация}
Для $I_n(D)=D^n(\ker D^{n+1})$:
\[
\boxed{I_n(D_\ell)=x^{a_\ell(n)}B,\qquad
a_\ell(n)=\ell n-\sum_{j=1}^{\ell}\floor{\frac n{2^j}}.}
\]
Доказательство даёт явный базис. Положим $v_j=r_{\ell-1-j}$ для $0\le j\le\ell$, $v_{-1}=p$. Тогда
\[
v_j^2=xv_{j+1}-v_{j-1}\quad(0\le j<\ell),\qquad \deg_{D_\ell}v_j=2^j.
\]
По моническим соотношениям нормальные мономы имеют вид
\[
s_n=\prod_{j=0}^{\ell-1}v_j^{\varepsilon_j}\,v_\ell^q,
\qquad n=\sum_{j=0}^{\ell-1}\varepsilon_j2^j+q2^\ell,
\quad\varepsilon_j\in\{0,1\}.
\]
Они образуют $B$-базис и имеют попарно разные $D$-степени. В выражении через $w$ старший коэффициент $v_j$ равен $x^{-(2^j-1)}$. Применение $D_\ell^n$ и двоичная запись $n$ дают указанную сумму целых частей; множитель $n!$ является единицей в $k$.

При $\ell=3$ показатели для $0\le n\le8$:
\[
0,\ 3,\ 5,\ 8,\ 9,\ 12,\ 14,\ 17,\ 17.
\]
При $\ell\ge2$ дифференцирование $D_\ell$ не сопряжено ни с одним неприводимым одношаговым переносом теоремы 8.1, даже при изменении базовой координаты. Действительно, единственный простой делитель plinth-идеала и его кратность $\ell$ сохраняются при любом сопряжении. Для одношагового класса показатель $I_2$ на этом простом был бы $\ell$ при $e=2$ или $2\ell$ при $e>2$, тогда как здесь он равен $2\ell-1$.

В частности, $D_\ell$ при $\ell\ge2$ \textbf{не триангулируемо}. Если бы в некоторых координатах оно имело вид $a(X)\partial_Y+b(X,Y)\partial_Z$ с $a\ne0$, оно принадлежало бы одношаговому классу с $h=a(X)Z-\int b(X,Y)\,dY$. Если первая координата имеет ненулевую постоянную производную, plinth равен всему ядру; случай $a=0$ при неприводимости также даёт постоянный перенос. Оба случая несовместимы с нашим plinth. Само существование таких нетриангулируемых примеров не объявляется новым; здесь дан явный способ распознать их двумя идеалами.

\subsection*{Ассоциированное градуированное кольцо}
Те же монические соотношения дают более полное описание:
\[
\boxed{\gr_{D_\ell}A\simeq
B[\xi_0,\ldots,\xi_\ell]/(\xi_j^2-x\xi_{j+1}:j<\ell),\quad
\deg\xi_j=2^j.}
\]
После редукции по $x$ нильрадикал $N=(\xi_0,\ldots,\xi_{\ell-1})$ удовлетворяет $N^\ell\ne0$, $N^{\ell+1}=0$, а $N/N^2$ свободен ранга $\ell$ над $k[p,\xi_\ell]$.

Показательный пример: $D_3$ и $xD_2$ имеют одинаковый plinth $(x^3)$, и показатели $I_n$ совпадают при $n\le7$. При $n=8$ они равны соответственно $17$ и $18$. Полная фильтрация видит различие, которое первые семь идеалов не обнаруживают.

Это иллюстрация дополнительной информации фильтрации, а не новая трудная
задача сопряжения: сами поля уже различаются неприводимостью.
При $\ell=2$ нетриангулируемость также известна из
\cite[Example 4.19]{DG}; новизна общего семейства и формул здесь
не устанавливается.

\subsection*{Связь с актуальным вопросом о сложности}
В журнальной статье Dasgupta--Gaifullin \cite[Theorem 4.12]{DG} universal level связан со сложностью биррационального представления формулой $\operatorname{ulev}(D)=c(\alpha)+1$. Их Question 4.21 относится к кандидату, построенному тремя чередующимися модификациями и перестановками.

Наша рекурсия позволяет исследовать соответствующий тип конструкции через точную фильтрацию. \textbf{Но равенство $c(\varphi)=3$ здесь не доказано}: для него требуется исключить все более короткие допустимые разложения, а не только канонический одношаговый класс. Возможность извлечь такую нижнюю оценку из $\gr_DA$ --- содержательная следующая задача.

\section{Почему относительный ML не различает наши кривые}\label{sec:relativeML}
Пусть $A=k[x_1,\ldots,x_n]$, $I\subsetneq A$ и $\dim(A/I)\le n-2$. Для любого $r\ge1$:
\[
\boxed{\bigcap_{\substack{0\ne D\in\operatorname{LND}(A)\\D(A)\subset I^r}}\ker D=k.}
\]
Более того, алгебра, порождённая объединением этих ядер, равна $A$.

Для каждого $i$ существует ненулевой $f_i\in I\cap k[x_1,\ldots,\widehat{x_i},\ldots,x_n]$: иначе $n-1$ образов координат были бы алгебраически независимы в кольце размерности $\le n-2$. Возьмём
\[
D_i=f_i^r\partial_{x_i},\qquad
\ker D_i=k[x_1,\ldots,\widehat{x_i},\ldots,x_n].
\]
Эти LND удовлетворяют $D_i(A)\subset I^r$, пересечение их ядер равно $k$, а объединение содержит все координаты. Это доказывает оба утверждения.

Следовательно, кандидат с более слабым условием $D(I)\subset I$ тоже всегда даёт $k$ для кривой в $\A^3$. Даже требование поточечно фиксировать произвольное конечное утолщение не помогает.

Для гармонических моделей можно взять, например,
\[
f_x=C_c(y)-C_b(z),\quad
f_y=C_c(x)-C_a(z),\quad
f_z=C_b(x)-C_a(y).
\]
Механизм реплик согласуется с более сильными результатами о гибкости дополнений, см. Flenner--Kaliman--Zaidenberg \cite{FKZ}. Поэтому здесь исправляется исследовательская стратегия, а не объявляется новый нетривиальный инвариант.


% This section is intended for inclusion in the consolidated Russian appendix.
% Bibliography keys used below: vdER, Shastri, Craighero, Masuda.
% All conclusions depend on the complete rectifier, not merely on a parameter lift.
\section{Какие именно именные гипотезы затрагивает выпрямление}

В этом разделе слово «опровержение» означает предъявленный полный
математический аргумент с точными сертификатами, представленный на проверку.
Внешнее рецензирование и литературный приоритет не утверждаются.

Положим $K=\mathbb Q(\sqrt{-3})$ и
\[
\gamma_T(u)=(f(u),g(u),h(u))
=(u^3-3u,u^4-4u^2,u^5-10u).
\]
В \cite[с.~1319--1320]{vdER} сформулированы три эквивалентные
гипотезы о конкретном объекте, связанном с этим вложением. Они
предполагают неполиномиальность соответственно ядра определённого LND,
факторкольца по его срезу и расширенного кольца Риса идеала кривой.
Последняя формулировка обозначена там как Asanuma's Conjecture.
Приведённый далее аргумент, применённый к удостоверенному выпрямителю
над $K$, даёт отрицание всех трёх формулировок над $\mathbb C$.

\begin{proposition}[Явные координаты связанного LND]
Пусть $k$ --- поле характеристики нуль,
$F=(P,Q,R_3)\in\operatorname{Aut}_k(\mathbb A^3)$,
$G=F^{-1}$, $F\circ\gamma=(u,0,0)$ и $0\ne b(t)\in k[t]$.
Для
\[
A=k[t,u,X_1,X_2,X_3],\qquad
D_b=b(t)\partial_u+\sum_i f_i'(u)\partial_{X_i}
\]
существует полиномиальная система координат
$A=k[t,p,q,r,s]$, в которой $D_b=\partial_s$.
\end{proposition}

\begin{proof}
Положим $V_i=f_i(u)-b(t)X_i$ и
\[
p=P(V),\qquad q=\frac{Q(V)}{b(t)},\qquad
r=\frac{R_3(V)}{b(t)},\qquad
s=\frac{u-P(V)}{b(t)}.
\]
Все числители точно делятся на $b(t)$, поскольку
$V\equiv\gamma(u)\pmod{b(t)}$. Обратная замена задаётся формулами
\[
u=p+b(t)s,\qquad
X_i=\frac{f_i(p+b(t)s)-G_i(p,b(t)q,b(t)r)}{b(t)}.
\]
Она также полиномиальна благодаря $G(p,0,0)=\gamma(p)$.
Обе композиции тождественны после обращения $b(t)$; инъективность
локализации доказывает их тождественность в исходном кольце. Наконец,
$D_b(V_i)=0$, откуда $D_b(p)=D_b(q)=D_b(r)=0$ и $D_b(s)=1$.
\end{proof}

Для модели Шастри и $b(t)=t$ получаем
\[
\ker D=K[t,p,q,r]\simeq K^{[4]}.
\]
После расширения скаляров это противоречит неполиномиальности,
утверждаемой в Conjecture~5.1 \cite{vdER}.

Исходный полином параметра Шастри равен
\[
F_0(x,y,z)=yz-x^3-5xy+2z-7x,
\]
а срез в Conjecture~5.2 равен
$\sigma=(u-F_0(V))/t$. Так как $D(\sigma)=D(s)=1$,
разность $\sigma-s$ принадлежит $B=\ker D$.
Следовательно, $A=B[\sigma]$; в частности,
\[
A/(\sigma-c)\simeq B\simeq K^{[4]}\qquad(c\in K).
\]
Это отрицание Conjecture~5.2 после перехода к $\mathbb C$.
При этом сам $F_0$ не является координатой $\mathbb A^3$:
его градиент исчезает в точках $(1,-2,5)$ и $(-1,-2,-5)$.
Некоординатность $F_0$ и координатность $\sigma$ совместимы,
поскольку речь идёт о разных полиномах в разных кольцах.

Наконец, при $R_0=K[x,y,z]$ и $I=I_{\gamma_T}$ имеем
$R_0=K[P,Q,R_3]$, $I=(Q,R_3)$, поэтому
\[
R_0[t,I/t]=K[t,P,Q/t,R_3/t]\simeq K^{[4]}.
\]
Здесь $R_0[t,I/t]\subset R_0[t,t^{-1}]$ --- именно расширенное
кольцо Риса, фигурирующее в Conjecture~5.3. Это отрицание третьей
формулировки, а не результат о произвольных кольцах Риса.

Таким образом, устранён один конкретный кандидат в трёх связанных
представлениях. Общая проблема сокращения из этого не решается.
Не решается и общая задача выпрямления $\mathbb A^1\hookrightarrow
\mathbb A^3$ либо общая гипотеза существования невыпрямимых
пространственных прямых.

\paragraph{Исторические различия.}
Предположение о комплексной невыпрямимости именно $\gamma_T$
приписано Шастри в \cite[с.~1319]{vdER}; предъявленный выпрямитель
его опровергает. Но вопросы Q.1 и Q.2 самой статьи
\cite[с.~12]{Shastri} уже решены Шастри: они относятся к
полиномиальным представлениям вещественных узлов, комплексирующимся
в замкнутые вложения, а не к их комплексной выпрямимости.
Специальные гипотезы Абъянкара, сформулированные через степени,
имели прежние опровержения, описанные в
\cite[введение, с.~115--116]{Craighero}.

\paragraph{Поля и действия.}
Полученные координаты определены над $K$ и затем над $\mathbb C$;
спуск на $\mathbb Q$ или $\mathbb R$ не доказан.
При $b=t$ естественное действие
$t\mapsto\lambda t$, $u\mapsto u$, $X_i\mapsto\lambda^{-1}X_i$
линейно в координатах ядра с весами $(1,0,-1,-1)$.
Это конкретная линеаризация. Утверждения о всех групповых действиях
из неё не следуют. В обратных связях с вложением существенна
градуировка, как и в работе \cite{Masuda}; её нельзя заменять
лишь абстрактной полиномиальностью ядра или равенством $\operatorname{ML}(B)=k$.

\subsection{Полный ML, относительный ML и фильтрация}
Для выпрямленной модели $B=k[t,p,q,r]$ четыре координатных
дифференцирования дают
\[
\ML(B)=k,\qquad \ML^*(B)=k,\qquad \operatorname{Dk}(B)=B.
\]
Здесь $\ML^*$ использует только LND со срезами, а $\operatorname{Dk}$
есть алгебра, порождённая ядрами всех ненулевых LND.
Для $D_b=\partial_s$ имеем $\pl(D_b)=B$,
$\ker D_b^{n+1}=\bigoplus_{j=0}^n Bs^j$ и
$D_b^n(\ker D_b^{n+1})=B$.
Это согласуется с большой минимальной степенью среза в исходных
координатах: речь идёт о разных характеристиках.

Для произвольного вложения, когда выпрямитель неизвестен, всё же верна
дихотомия
\[
\ML(B_b)\in\{k,k[t]\},\qquad B_b=\ker D_b.
\]
Действительно, $\partial_{X_i}$ коммутируют с $D_b$ и ограничиваются
на $B_b$. Пересечение их ядер там равно
$B_b\cap k[t,u]=k[t]$. Если непостоянный $a(t)$ лежит в полном ML,
то для любого LND $\Delta$ на $B_b$ равенство
$0=\Delta a(t)=a'(t)\Delta t$ в области характеристики нуль
влечёт $\Delta t=0$. Тогда весь $k[t]$ лежит в ML.

Поэтому $\ML(B_b)=k[t]$ было бы препятствием выпрямлению.
Однако равенство $\ML(B_b)=k$ не доказывает ни полиномиальности $B_b$,
ни выпрямимости. Если с самого начала разрешать только LND,
фиксирующие $t$, относительный ML равен $k[t]$ даже в полиномиальном
случае. Для новых ядер $H(3,7,8)$ и $H(4,5,7)$ выбор ветви дихотомии
пока не выполнен.

\subsection{Что надо сохранить для обратного перехода}
Пусть $R=k[x,y,z]$, $I=I_\gamma$,
$B_\mu=R[t,I/t^\mu]$, $\mu\ge1$, $\wt(t)=1$, $\wt(R)=0$.
Вложение выпрямимо тогда и только тогда, когда существуют однородные
координаты
\[
B_\mu=k[t,p,q,r],\qquad
\wt(t,p,q,r)=(1,0,-\mu,-\mu)
\]
с \emph{тем же выделенным элементом $t$}. Обратное направление следует
из равенств
\[
\begin{gathered}
(B_\mu)_0=R=k[p,t^\mu q,t^\mu r],\\
I=t^\mu(B_\mu)_{-\mu}=R\cap tB_\mu=(t^\mu q,t^\mu r).
\end{gathered}
\]
Ограничение $p$ на кривую является аффинным параметром; его можно
нормировать до $u$.

Эквивалентная формулировка: существует однородное LND $E$ степени $-1$
на $B_\mu$ с $E(t)=1$. Прямое направление даёт $\partial_t$.
Обратно, теорема о срезе и конечная проекция
\[
\rho(b)=\sum_{j\ge0}\frac{(-t)^j}{j!}E^j(b)
\]
дают $B_\mu=(\ker E)[t]$ и градуированный изоморфизм
$\ker E\simeq B_\mu/(t)$. Последний слой полиномиален с весами
$(0,-\mu,-\mu)$ по построению нормального базиса в разделе~\ref{sec:formal}.
Значит, поднимаются требуемые однородные координаты.

Для LND $E$, без предположения однородности, условие $E(t)=1$ ещё даёт абстрактную полиномиальность
$B_\mu$, но приведённое доказательство уже не восстанавливает вложение.
Для LND $E$ с одним лишь $E(t)\ne0$ следует только $\ML(B_\mu)=k$.
Нельзя без доказательства брать внутреннюю однородную компоненту
произвольного LND и объявлять её локально нильпотентной.

\paragraph{Непосредственная связь с плоской теоремой.}
Если идеал $I$ содержит полиномиальную координату, кривую можно
поместить в координатную плоскость и применить теорему
Абъянкара--Моха--Судзуки. Если координатой является полином параметра
$P\circ\gamma=u$, то при $R=k[P,U,V]$ нужный выпрямитель сразу равен
$(P,U-U(\gamma(P)),V-V(\gamma(P)))$.
Обычного, не обязательно координатного полинома параметра для обоих
рассуждений недостаточно.

\section{Нормальное поле Эйлера: полный конструктивный критерий}
Пусть $R=k[Y_1,\ldots,Y_n]$, $\gamma:\A^1\hookrightarrow\A^n$ --- замкнутое вложение,
$I=I_\gamma$, $B=R[t,I/t]$. Для дифференцирования $\delta$ рассматриваем
условия
\begin{equation}\label{eq:normal}
\delta(R)\subset I,\qquad (\delta-\mathrm{id})(I)\subset I^2.
\end{equation}
Они означают исчезновение поля на кривой и единичное действие на
линейных нормальных функциях. Тогда
\[
E=\partial_t+t^{-1}\delta
\]
сохраняет $B$, поскольку $E(Y_i)=\delta Y_i/t\in B$ и
$E(F/t)=(\delta F-F)/t^2\in B$ для $F\in I$.
Локальная конечность $\delta$ означает конечномерность
$\Span_k\{\delta^jF:j\ge0\}$ для каждого $F\in R$.

\begin{theorem}\label{thm:euler}
Для заданного $\delta$, удовлетворяющего \eqref{eq:normal}, эквивалентны:
\begin{enumerate}
\item $E$ является LND на $B$;
\item $\delta$ локально конечно на $R$;
\item существуют координаты $R=k[p,q_1,\ldots,q_{n-1}]$ с
$p\gamma=u$, $I=(q_1,\ldots,q_{n-1})$ и
$\delta=\sum_jq_j\partial_{q_j}$.
\end{enumerate}
Следовательно, существование такого локально конечного поля
равносильно выпрямимости $\gamma$ над $k$.
\end{theorem}

\begin{proof}
В локализации непосредственно по индукции получаем
\begin{equation}\label{eq:falling}
E^mF=t^{-m}\prod_{j=0}^{m-1}(\delta-j)F,
\qquad
E^m(F/t)=t^{-m-1}\prod_{j=1}^{m}(\delta-j)F.
\end{equation}
Если $E$ локально нильпотентно, каждый $F$ аннулируется ненулевым
полиномом от $\delta$, что даёт локальную конечность.

Правило Лейбница и \eqref{eq:normal} дают
\[
(\delta-m)I^m\subset I^{m+1}\quad(m\ge0).
\]
Предположим теперь, что $\delta$ локально конечно.
Пусть $V$ --- конечномерное $\delta$-стабильное пространство,
$M$ --- максимальная обычная степень его элементов.
Так как $I\subset\mathfrak m_{\gamma(0)}$, имеем $V\cap I^{M+1}=0$.
Поэтому $\delta(\delta-1)\cdots(\delta-M)$ аннулирует $V$.
Его корни различны и лежат в $k$: полупростота и неотрицательные
целочисленные веса следуют автоматически, без перехода к
алгебраическому замыканию. Получаем
\[
R=\bigoplus_{d\ge0}R_d,\qquad \delta|_{R_d}=d.
\]
Для ненулевого собственного вектора точный $I$-адический порядок
равен его весу, поскольку на $I^m/I^{m+1}$ действие равно $m$.
Идеалы $I^m$ градуированы: спектральные проекторы являются
полиномами от $\delta$. Отсюда
\[
I^m=\bigoplus_{d\ge m}R_d,\quad R_0\simeq R/I\simeq k[u],\quad
R_1\simeq I/I^2,\quad R=R_0[R_1].
\]
Последнее равенство следует индукцией по весу: при $d\ge2$ имеем
$R_d\subset I^2$, и каждый однородный множитель в произведении
имеет меньший положительный вес.

Для свободности $I/I^2$ не требуется гомологическая алгебра.
Выберем $P\gamma=u$, положим $v=\gamma'$, $a=\nabla P(\gamma)$.
Тогда $av=1$, а модуль строк
$L=\{\ell\in k[u]^n:\ell v=0\}$ свободен ранга $n-1$
над PID $k[u]$. Градиент задаёт изоморфизм
\[
I/I^2\longrightarrow L,\qquad F\longmapsto\nabla F(\gamma).
\]
Явный подъём строки есть
\[
F_\ell(Y)=\sum_i\ell_i(P(Y))\bigl(Y_i-f_i(P(Y))\bigr).
\]
Инъективность доказывается Тейлором: функция с нулевым значением
и нулевым первым градиентом на кривой принадлежит $I^2$.
Выбираем инвариантный подъём $p$ параметра $u$ и базис $q_j$ модуля
$R_1$ над $R_0=k[p]$. Тогда $R=k[p,q_1,\ldots,q_{n-1}]$;
алгебраическая независимость следует из степени трансцендентности $n$.
Получена нормальная форма $\delta$.
Наконец, $B=k[t,p,q_1/t,\ldots,q_{n-1}/t]$, и $E=\partial_t$
в этих координатах. Это доказывает все импликации.
\end{proof}

\subsection{Конечная проверка и восстановление автоморфизма}
Достаточно предъявить точную стабильность
\[
V_N=\Span_k\{1,\delta^jY_i:0\le j\le N, 1\le i\le n\},\qquad
\delta^{N+1}Y_i\in V_N.
\]
Пространства произведений элементов $V_N$ стабильны по Лейбницу,
поэтому $\delta$ локально конечно на всём $R$.
Если степени элементов $V_N$ не превосходят $M$, спектральная
проекция полинома $F$ на вес $j$ вычисляется как
\[
\pr_j(F)=\prod_{\substack{0\le i\le L\\i\ne j}}
\frac{\delta-i}{j-i}F,\qquad L\ge M\deg F.
\]
Таким образом,
$p=\pr_0(P)$, $q_j=\pr_1(F_{\ell_j})$ --- искомые координаты.
Обратное отображение восстанавливается по весовым компонентам
$Y_i$ и конечным поперечным производным. Подробная формула сохранена
в \file{euler/theory.md}. Универсальной априорной границы $N$ здесь нет.

Положительный контроль восстанавливает нелинейные координаты
\[
r=z-x^2,\qquad q=y-r^2,\qquad p=x-q^2.
\]
Соответствующее пространство $V_4$ имеет размерность $9$ и точно
замыкается. Это контроль реализации, а не замена общего доказательства.

\subsection{Необходимые условия, которые ещё недостаточны}
Из нормальной формы следуют
\[
\Div\delta=n-1,\qquad (\delta Y_1,\ldots,\delta Y_n)=I.
\]
Равенство идеалов описывает всю неподвижную схему, а не только
множество точек. Постоянство дивергенции LF-поля является классическим
фактом \cite{Nowicki}; приоритет этого факта не заявляется.

Даже оба условия вместе с \eqref{eq:normal} не достаточны.
Для прямой $I=(y,z)$ поле
\[
\delta=yz\partial_x+y\partial_y+(z+y^2x)\partial_z
\]
проходит все эти тесты. Однако $\delta^{2m}x$ содержит
$xy^{3m}$ с положительным коэффициентом. Он получается повторением
пути $x\mapsto yz\mapsto y^3x$; все коэффициенты неотрицательны,
сокращений нет. Степени неограниченны, так что поле не LF.
Нельзя смешивать LF нормального поля с LND: в правильной нормальной
форме $\delta^m q_j=q_j$, а локально нильпотентным является $E$.

\section{Полная граница наивного нормального поля}
Для любого полинома параметра $P\gamma=u$ положим
\[
\delta_P=\sum_i\bigl(Y_i-f_i(P)\bigr)\partial_{Y_i}.
\]
Оно всегда удовлетворяет \eqref{eq:normal}, но имеет следующий
существенный предел применимости.

\begin{theorem}
Существует $P\gamma=u$, для которого $\delta_P$ локально конечно,
тогда и только тогда, когда существует аффинно-линейный $L$ с
$L\gamma=u$.
\end{theorem}
\begin{proof}
Если $P=L$ аффинно-линейно, то $\delta_P(P)=0$ и
$\delta_P^2Y_i=\delta_PY_i$, что даёт LF на образующих.
Обратно, необходимое равенство дивергенции имеет вид
\[
\sum_i f_i'(P)\partial_iP=1.
\]
Аффинной заменой координат приведём ненулевые положительные степени
$f_i$ к попарно различным; остальные компоненты можно сделать нулевыми.
Вид $\delta_P$ при этой замене сохраняется. Если линейного параметра
нет, все положительные степени не меньше $2$.
Пусть $d=\deg P$, а $j$ соответствует наибольшей степени $f_j$
с $\partial_jP\ne0$. Степень этого слагаемого не меньше
$(\deg f_j-1)d$, тогда как каждое предыдущее имеет степень не больше
$(\deg f_i)d-1\le(\deg f_j-1)d-1$.
Старший член не сокращается, и сумма не равна $1$.
Если все активные ненулевые компоненты отсутствуют, сумма равна нулю,
что также невозможно.
\end{proof}

Для гармонических моделей с различными индексами $a,b,c\ge2$
линейного параметра нет. Поэтому \emph{ни один} выбор $P$, какой бы
ни была его степень, не сделает это поле LF. Запрет относится также
к уже выпрямленным $H(3,4,5)$ и $H(3,5,7)$ и потому не является
препятствием выпрямимости самой кривой.

\section{Общие нормальные поля: точные пространства кандидатов}
Ищем все $\delta=a_1\partial_x+a_2\partial_y+a_3\partial_z$ с
$\deg a_i\le d$, условиями \eqref{eq:normal} и $\Div\delta=2$.
Каждый $a_i$ раскладывается по полному базису $I\cap R_{\le d}$.
При $v=\gamma'$ второе нормальное условие эквивалентно тому, что
столбцы $J_a(\gamma)-\mathrm{Id}$ кратны унимодулярному $v$.
Получается рациональная линейная система $Ac=b$, охватывающая
все поля данной степени.

\begin{center}\small
\begin{tabular}{@{}lrrrrl@{}}\toprule
Модель & $d$ & $\dim I_{\le d}$ & Неизвестных & $\operatorname{rank}A$ & Итог\\\midrule
$H(3,7,8)$ &7&63&189&189&несовместна\\
$H(3,7,8)$ &8&100&300&276&семья $24$\\
$H(4,5,7)$ &6&42&126&---&несовместна\\
$H(4,5,7)$ &7&71&213&194&семья $19$\\\bottomrule
\end{tabular}
\end{center}
Несовместность удостоверена точным рациональным вектором
$\lambda^{\mathsf T}A=0$, $\lambda^{\mathsf T}b=1$; поэтому вывод
сохраняется при любом расширении поля характеристики нуль.
Для совместных систем предъявлены рациональная точка и полный базис
направлений. Точные нулевые векторы дают верхние оценки ранга;
совпадающие модульные нижние оценки удостоверяют полноту.

Таким образом, первая степень совместности \emph{нормальных условий
вместе с глобальной дивергенцией $2$} равна $8$ и $7$.
Это не минимальные степени LF-полей или выпрямителей.
Для $H(4,5,7)$ степень $6$ невозможна уже без условия дивергенции.
Для $H(3,7,8)$ нормальные условия в степени $7$ сами по себе имеют
семейство размерности $11$; именно глобальная дивергенция исключает его.

\subsection{Два выбранных представителя строго исключены}
В семье $H(3,7,8)$ выбранный рациональный представитель сохраняет
$(x,y)$ и на факторе $k[z]$ действует как $\psi(z)\partial_z$,
$\deg\psi=6$. В семье $H(4,5,7)$ выбранный представитель сохраняет
$(y,z)$ и индуцирует $\varphi(x)\partial_x$, $\deg\varphi=7$.
Для полинома $a$ степени $e\ge2$ в характеристике нуль
\[
\deg((a\partial_v)^m v)=1+(e-1)m,\qquad
\operatorname{lc}((a\partial_v)^m v)
=\operatorname{lc}(a)^m\prod_{j=0}^{m-1}(1+(e-1)j)\ne0.
\]
Поэтому оба представителя не LF, а соответствующие $E$ не LND.
Полные коэффициенты и отдельные проверки приведены в JSON-сертификатах.
Оси не сохраняются всем семейством: соответственно $13$ из $24$ и
$11$ из $19$ выбранных базисных направлений их нарушают.
Эти числа относятся к выбранному базису; это не размерности
исключённых подсемейств. Общего запрета на все комплексные параметры нет.

\subsection{Вещественная топология и сопряжение}
По \cite[с.~2, рис.~3]{KP}, $H(3,7,8)$ имеет тип $T(2,5)$
с точностью до зеркала, а $H(4,5,7)$ --- $5_2$.
Следовательно, вещественных полиномиальных выпрямителей у этих моделей
нет. По теореме~\ref{thm:euler} \emph{любое} хорошее LF-поле для них
должно иметь невещественные коэффициенты. Это известное топологическое
ограничение поля определения, а не комплексное препятствие.

Если хорошее комплексное $\delta$ коммутировало бы с $\bar\delta$,
их среднее было бы вещественным LF-полем с теми же нормальными
условиями. Коммутирующие LF-операторы имеют конечномерные совместные
циклические пространства, поэтому среднее действительно LF.
Получилось бы запрещённое вещественное выпрямление. Аналогично,
коммутирующие $E$ и $\bar E$ дали бы вещественное LND со срезом $t$.
Следовательно, в этих случаях оба коммутатора ненулевые. При этом
\[
[E,\bar E]=t^{-2}\bigl([\delta,\bar\delta]+\delta-\bar\delta\bigr),
\]
так что два условия коммутирования нельзя отождествлять.

\section{Что осталось гипотетическим и как продолжать}

Установленные здесь утверждения решают конкретные задачи для конкретных
моделей и классов дифференцирований. Следующая таблица фиксирует задачи,
которые нельзя записывать в список решённых.

\begin{longtable}{L{0.30\textwidth}L{0.64\textwidth}}
\toprule Задача & Точный статус и следующий проверяемый шаг \\ \midrule
\endhead
Все гармонические вложения & Даны два бесконечных семейства и конечные
проверки. Нужен общий переход между классами вычетов либо новое
конструктивное выпрямление. Ограниченный перебор не является классификацией.\\[6pt]
Комплексное выпрямление $H(3,7,8)$ и $H(4,5,7)$ & Не построено.
Можно искать локально конечное поле в найденных аффинных пространствах
кандидатов. Отсутствие такого поля в этих степенях, если будет доказано,
не исключит поля большей степени.\\[6pt]
Полный ML новых ядер & Пока имеется дихотомия $k$ или $k[t]$.
Для результата $k$ достаточно одного LND, двигающего $t$.
Для результата $k[t]$ надо доказать, что каждый LND фиксирует $t$;
ограниченный поиск этого не доказывает.\\[6pt]
Однородный $E$ со срезом $t$ & Условия $E(t)=1$ и сохранения ядра
ещё не дают LND. Обнуление итераций на образующих является точным
положительным сертификатом; неудача конечного числа итераций не
доказывает отсутствия подходящего поля.\\[6pt]
Формальные координаты & Построены все конечные утолщения и пополнение.
Не доказано, что координаты обрываются в полиномы: это отдельная
глобальная задача.\\[6pt]
Question 4.21 Dasgupta--Gaifullin & Формулы фильтраций получены,
но минимальность трёхшагового биррационального представления не
установлена. Нужно исключить все допустимые более короткие представления.\\[6pt]
Минимальные степени выпрямителей & Степени предъявленных автоморфизмов
являются верхними оценками. Минимумы степеней срезов не заменяют эту задачу.\\[6pt]
Вещественный спуск & Вещественное выпрямление рассматриваемых нетривиальных
узлов невозможно. Полиномиальность соответствующего вещественного
ядра без сохранения выделенного $t$ этим сама по себе не исключается.\\[6pt]
Общие задачи сокращения и линеаризации & Отвергнут один предложенный
кандидат; общие задачи и все возможные действия не исследованы полностью.\\
\bottomrule
\end{longtable}

\subsection*{Почему поиск LND со срезом мотивирован}
В кольце $B_\mu$ элемент $t$ выделен самим вложением, а специальный слой
$t=0$ уже имеет полиномиальное описание. Если удастся предъявить
однородное дифференцирование $E$ степени $-1$ с $E(t)=1$ и доказать его
локальную нильпотентность, формула среза даст
\[
 \rho_E(f)=\sum_{j\ge0}\frac{(-t)^j}{j!}E^j(f),\qquad
 B_\mu=(\ker E)[t].
\]
Сумма конечна именно благодаря LND. При однородности эта конструкция
возвращает координаты исходного вложения, как доказано в \S11.
Поэтому дополнительное требование LND отвечает точной геометрической
цели. Оно не является технической деталью, автоматически следующей
из линейных уравнений на коэффициенты.

Для конечного набора образующих $b_1,\ldots,b_m$ достаточно предъявить
целые $N_i$ и тождества $E^{N_i}(b_i)=0$. Формула Лейбница тогда даёт
локальную нильпотентность на всей алгебре. Для поиска нормального
локально конечного поля альтернативный положительный сертификат ---
конечномерное устойчивое пространство, содержащее образующие, из \S12.
Оба способа дают проверяемое доказательство в случае успеха,
но здесь не доказана общая граница времени или степени для их поиска.

\subsection*{Роль групп Кремоны}
Биррациональные замены вида
\[
 (x,y,z)\longmapsto (x,y,m(x)z-q(x,y))
\]
дают эффективный язык переноса полей и координат после локализации.
Возвращение к полиномиальному кольцу требует точной делимости и проверки
центра модификации. Установленный запрет для естественного центра
семейства $H(3,3n+1,3n+2)$ показывает, что формального рационального
обратного отображения недостаточно. Практическая программа состоит
в поиске других центров, в вычислении образных идеалов и в сравнении
различных факторизаций. Общая классификация подгрупп Кремоны здесь
не заявляется.

\subsection*{Связь с теоремой Абъянкара--Моха}
Прямое применение плоской теоремы возможно после предъявления
координатной поверхности, содержащей кривую: на этой поверхности
вложенная прямая становится плоским вложением, к которому применяется
\cite{AM,Suzuki}. Найти такую поверхность в пространстве ---
содержательная недостающая часть. Эквивалентные критерии через
координатный параметр, градуированное LND и нормальное локально конечное
поле позволяют атаковать именно эту часть разными методами.
Не следует подменять её уже известной плоской теоремой.

\subsection*{Порядок чтения}
Для введения в LND полезно начать с \emph{First Principles} и
\emph{Further Properties of LNDs} книги Freudenburg \cite{Freudenburg},
затем перейти к главам \emph{Dimension Three}, \emph{Algorithms},
\emph{Makar-Limanov and Derksen Invariants} и
\emph{Slices, Embeddings and Cancellation}. Для настоящего проекта
следующий круг чтения --- \cite{vdER} (точные гипотезы и кольцо ядра),
\cite{Masuda} (эквивариантная структура), \cite{EVM} (относительные
координаты), \cite{DG} (итерации LND и сложность),
\cite{FF,KP} (гармонические кривые и узлы).

\section{Воспроизводимость и состав приложения}

Все вычислительные файлы помещены в каталог \file{computations/}.
Комментарии в поставляемом коде даны на английском языке. Исходные
версии сохранены в архиве происхождения; отдельно проверяется, что
редактирование комментариев и документационных строк Python не
меняет абстрактное синтаксическое дерево исполняемой части.
Это полезный контроль упаковки, но не доказательство правильности
алгоритма.

\begin{longtable}{L{0.29\textwidth}L{0.65\textwidth}}
\toprule Каталог & Назначение \\ \midrule
\endhead
\file{baseline/} & Исходные выпрямители, точные проверки взаимно
обратных отображений, напечатанных компонент и исходных сертификатов.\\[5pt]
\file{legacy_appendix/} & Вычисления опубликованного аппендикса
от 26 октября, включая локальные зависимости для автономного запуска.\\[5pt]
\file{continuation/} & Четырнадцать проверок продолжения: минимумы
срезов, степени кратности, формальные координаты, ограничения
переноса, фильтрации LND и относительный ML.\\[5pt]
\file{euler/} & Нормальные поля, ранги линейных систем, двойственные
сертификаты несовместности и осевые препятствия для выбранных полей.\\[5pt]
\file{audit_checks/} & Проверка отождествления кандидата из
Conjectures 5.1--5.3 с исходной моделью и исходным срезом.\\[5pt]
\file{verification_logs/} & Стандартный вывод, ошибки и время
каждой программы повторного прогона.\\[5pt]
\file{provenance/} & Исходные версии, контрольные суммы и сведения
о происхождении материалов.\\
\bottomrule
\end{longtable}

Точка входа --- \file{computations/run_all_checks.py}.
Запуск следует выполнять из каталога \file{computations/}:
\begin{quote}\ttfamily\small
python -m pip install -r requirements.txt\\
python run\_all\_checks.py
\end{quote}
Свежий запуск от 28 октября 2026 года завершился успешно: \textbf{35 из 35
программ}. Все основные алгебраические проверки выполняются точно, без замены
тождеств численным приближением. Итоговый машинный протокол
\file{verification_summary.json} фиксирует результаты и команды.
Число программ не равно числу теорем: одна программа может проверять
много тождеств, а бесконечные семейства требуют приведённых выше
математических доказательств.

Код и исторические протоколы GAP/Macaulay2 включены для независимой
проверки, но эти системы в текущем окружении не запускались.
Их старые успешные протоколы явно отделены от нового Python-прогона.
Длительные поисковые эксперименты сохраняются как исходники и данные;
они не подменяют краткие проверяющие программы.

\subsection*{Совместимость с репозиторием}
Базовая версия проверена по коммиту
\file{b954f3609e900761fddc9c8ff03d8301e9604cab}.
Все 75 файлов опубликованного аппендикса совпали с Git-объектами этой
версии. Обновление добавляет только отдельный каталог
\file{updates/2026-10-28/} и указатель на него внутри прежнего
аппендикса. Основной препринт и прежние сертификаты не заменяются.

Содержательно уточняются две прежние исследовательские перспективы.
Во-первых, относительный ML старого аппендикса предлагался условно;
\S10 показывает, что ожидаемая нетривиальность при данном определении
невозможна. Во-вторых, новая граница степени $\ge5$ для $H(4,5,7)$
относится к координатам, сохраняющим одну из исходных переменных.
Общий поиск абсолютной координаты со степени $4$ остаётся обоснованным.
Такое уточнение не требует изменения старого архивного текста,
но должно сопровождать его дальнейшее цитирование.

\subsection*{Проверка читателем}
Сначала следует проверить формулировку конкретной модели и поле
коэффициентов; затем запустить сертификаты выпрямителя; после этого
проверить конструктивный переход к ядру и срезу. Для новых общих
утверждений основным объектом проверки остаётся математическое
доказательство. При обнаружении ошибки необходимо указать модель,
утверждение, команду запуска и полный протокол. Автор приглашает
читателей к такой проверке и к развитию исследовательской программы.

\begin{thebibliography}{99}
\small\setlength{\itemsep}{6pt}
\bibitem{Repository} I. Dorofeev.
\emph{Polynomial Rectification of the Shastri Trefoil and a Harmonic
Figure-Eight}. Рабочий препринт и вычислительные сертификаты.
\href{https://github.com/qft98romanov-hash/abhyankar-moh-trefoil/tree/b954f3609e900761fddc9c8ff03d8301e9604cab}{Репозиторий, проверенная версия}.

\bibitem{AM} S. S. Abhyankar, T. T. Moh.
\emph{Embeddings of the line in the plane}.
J. reine angew. Math. \textbf{276} (1975), 148--166.
\href{https://doi.org/10.1515/crll.1975.276.148}{DOI: 10.1515/crll.1975.276.148}.

\bibitem{Suzuki} M. Suzuki.
\emph{Propriétés topologiques des polynômes de deux variables complexes,
et automorphismes algébriques de l'espace $\C^2$}.
J. Math. Soc. Japan \textbf{26} (1974), 241--257.
\href{https://doi.org/10.2969/jmsj/02620241}{DOI: 10.2969/jmsj/02620241}.

\bibitem{Shastri} A. R. Shastri.
\emph{Polynomial representations of knots}.
Tôhoku Math. J. (2) \textbf{44} (1992), 11--17.
\href{https://doi.org/10.2748/tmj/1178227371}{DOI: 10.2748/tmj/1178227371};
\href{https://www.jstage.jst.go.jp/article/tmj1949/44/1/44_1_11/_pdf}{открытый PDF}.
Вопросы Q.1--Q.2 и Theorem 1, с.~12; Example 1, с.~15.

\bibitem{Craighero} P. C. Craighero.
\emph{About Abhyankar's conjectures on space lines}.
Rend. Sem. Mat. Univ. Padova \textbf{74} (1985), 115--122.
\href{https://www.numdam.org/item/RSMUP_1985__74__115_0.pdf}{Открытый PDF}.
Исторические формулировки: введение, с.~115--116.

\bibitem{vdER} A. van den Essen, P. van Rossum.
\emph{Triangular derivations related to problems on affine $n$-space}.
Proc. Amer. Math. Soc. \textbf{130} (2002), no.~5, 1311--1322.
\href{https://doi.org/10.1090/S0002-9939-01-06212-8}{DOI: 10.1090/S0002-9939-01-06212-8}.
Conjectures 5.1--5.3, Proposition 5.4, Corollary 5.5, с.~1319--1321.

\bibitem{Masuda} K. Masuda.
\emph{Equivariant embeddings and semi-invariant locally nilpotent
derivations}. Авторская версия.
\href{https://sci-tech.ksc.kwansei.ac.jp/~kayo/cmn/papers/eembed.pdf}{Открытый PDF}.
Theorems 1.1--1.2; градуировка является существенной частью постановки.

\bibitem{EVM} E. Edo, A. van den Essen, S. Maubach.
\emph{A note on $k[z]$-automorphisms in two variables}.
J. Pure Appl. Algebra \textbf{213} (2009), 1197--1200.
\href{https://doi.org/10.1016/j.jpaa.2008.11.040}{DOI: 10.1016/j.jpaa.2008.11.040};
\href{https://arxiv.org/abs/0809.0767}{arXiv:0809.0767}.
Proposition 3.1 в открытой версии.

\bibitem{FF} G. Freudenburg, J. Freudenburg.
\emph{Curves defined by Chebyshev polynomials} (2009).
\href{https://arxiv.org/abs/0902.3440}{arXiv:0902.3440}.

\bibitem{KP} P.-V. Koseleff, D. Pecker.
\emph{Harmonic knots}.
J. Knot Theory Ramifications \textbf{25} (2016), 1650074.
\href{https://doi.org/10.1142/S0218216516500747}{DOI: 10.1142/S0218216516500747};
\href{https://arxiv.org/abs/1203.4376}{arXiv:1203.4376}.
Семейство $T(2,2n+1)$ и рисунок 3 для $H(4,5,7)$.

\bibitem{Chau} N. V. Chau.
\emph{A simple proof of Jung's theorem on polynomial automorphisms
of $\C^2$} (2004).
\href{https://arxiv.org/abs/math/0408077}{arXiv:math/0408077}.

\bibitem{DG} N. Dasgupta, S. A. Gaifullin.
\emph{On locally nilpotent derivations of polynomial algebra in three
variables}. Sbornik: Mathematics \textbf{216}:4 (2025), 456--484.
\href{https://doi.org/10.4213/sm10094e}{DOI: 10.4213/sm10094e}.
Русская версия: Математический сборник \textbf{216}:4 (2025), 3--34,
\href{https://doi.org/10.4213/sm10094}{DOI: 10.4213/sm10094}.
Theorem 4.12, Example 4.19, Remark 4.20, Question 4.21.

\bibitem{FKZ} H. Flenner, S. Kaliman, M. Zaidenberg.
\emph{A Gromov--Winkelmann type theorem for flexible varieties}.
J. Eur. Math. Soc. \textbf{18} (2016), 2483--2510.
\href{https://doi.org/10.4171/JEMS/646}{DOI: 10.4171/JEMS/646};
\href{https://ems.press/content/serial-article-files/32154}{открытый PDF}.

\bibitem{Nowicki} A. Nowicki.
\emph{The divergence of polynomial derivations}.
Авторский текст от 29 декабря 1993 года, Theorem 3.3.
\href{https://www.researchgate.net/publication/242147525_The_divergence_of_polynomial_derivations}{Открытая авторская копия}.

\bibitem{Freudenburg} G. Freudenburg.
\emph{Algebraic Theory of Locally Nilpotent Derivations}.
2nd enlarged edition, Encyclopaedia of Mathematical Sciences,
Springer, 2017.
\href{https://doi.org/10.1007/978-3-662-55350-3}{DOI: 10.1007/978-3-662-55350-3}.
\end{thebibliography}

\end{document}
